% !TeX program = xelatex
% ============================================================
% report/report.tex — course project report, 1st year of the HSE MSc
% programme «System and Software Engineering».
%
% Structure per the programme's course project guidelines: title page
% (App. 4) → Abstract → Contents → Definitions and abbreviations →
% Introduction → Main part → Conclusion → References → Appendices.
% The main part depends on the work format chosen at project creation:
%   project-research (§5.1) — analysis and problem statement, design,
%     implementation and testing; technical documentation as an appendix or
%     a link to the repository (§5.2);
%   academic (§5.6) — 3–5 chapters: review and problem statement,
%     theoretical results, implementation and/or experiment, analysis.
% Mentor seminar: chapter 1 carries the problem statement and all proposal
% sections (aim, tasks, methods, tools).
%
% Replace the yellow \hseFill{…} with your text.
% ============================================================
\documentclass[12pt,a4paper]{extreport}
% !TeX program = xelatex
% ============================================================
% report/preamble.tex — преамбула Отчёта о курсовом проекте магистратуры
% СПИ ФКН НИУ ВШЭ.
%
% Методичка ОП СПИ §5.1/§5.6: структура отчёта о НИР по ГОСТ 7.32 —
% поля 30/15/20/20, Times New Roman 12 пт, интервал 1,5, абзац 1,25 см,
% сквозная нумерация рисунков и таблиц. Чистый документ без штампа и
% рамки. Язык разделов — язык работы (project.lang).
% ============================================================
\newcommand{\hseDefaultMono}{Consolas}
\input{../common/fonts}

((* if project.lang == 'en' *))
\usepackage[russian,english]{babel}
\addto\captionsenglish{\renewcommand{\bibname}{References}}
((* else *))
\usepackage[english,russian]{babel}
\addto\captionsrussian{\renewcommand{\bibname}{Список использованных источников}}
((* endif *))


% Поля страницы — по настройке проекта report_margins (Настройки → Поля
% страницы отчёта). По умолчанию ГОСТ 7.32-2017 (30/15/20/20); вариант ВШЭ —
% 25/10/20/20. Значение подставляется при создании проекта.
((* if 'ВШЭ' in (project.meta.report_margins | default('')) *))
\usepackage[a4paper,left=25mm,right=10mm,top=20mm,bottom=20mm]{geometry}
((* else *))
\usepackage[a4paper,left=30mm,right=15mm,top=20mm,bottom=20mm]{geometry}
((* endif *))
\usepackage{setspace}
\onehalfspacing
\usepackage{indentfirst}
\setlength{\parindent}{1.25cm}
\setlength{\parskip}{6pt}

\usepackage{microtype}
\usepackage{fvextra}
\usepackage{csquotes}
\usepackage{amsmath}
\usepackage{amssymb}
\usepackage{etoolbox}
\usepackage{graphicx}
\usepackage{letltxmacro}
\LetLtxMacro{\oldincludegraphics}{\includegraphics}
\renewcommand{\includegraphics}[2][]{%
  \IfFileExists{#2}{%
    \oldincludegraphics[#1]{#2}%
  }{%
    \IfFileExists{#2.png}{\oldincludegraphics[#1]{#2.png}}{%
      \IfFileExists{#2.jpg}{\oldincludegraphics[#1]{#2.jpg}}{%
        \fcolorbox{red}{white}{%
          \begin{minipage}{0.5\textwidth}\centering
            \textcolor{red}{\textbf{Файл не найден:}} \\ \texttt{\detokenize{#2}}
          \end{minipage}%
        }%
      }%
    }%
  }%
}
\usepackage{xcolor}
\usepackage{hyperref}
\usepackage{lastpage}
\usepackage{float}
\usepackage{array}
\usepackage{booktabs}
\usepackage{longtable}
\usepackage{tabularx}
\usepackage{adjustbox}
\usepackage{xurl}
\urlstyle{same}
\usepackage{hyphenat}
\usepackage{caption}
% Ссылки [n] печатаются в строку (ГОСТ 7.0.5-2008), не надстрочно.
\usepackage{cite}
\DeclareCaptionLabelSeparator{emdash}{~---~}
\captionsetup[table]{position=top,labelsep=emdash,justification=raggedright,singlelinecheck=false}
\captionsetup[figure]{position=bottom,labelsep=emdash}

% ── Листинги кода ───────────────────────────────────────────
\usepackage{listings}
((* if project.lang != 'en' *))\renewcommand{\lstlistingname}{Листинг}((* endif *))
\definecolor{codegray}{rgb}{0.5,0.5,0.5}
\definecolor{backcolour}{rgb}{0.95,0.95,0.92}
\lstset{
  basicstyle=\footnotesize\ttfamily,
  backgroundcolor=\color{backcolour},
  commentstyle=\color{codegray}\itshape,
  numberstyle=\tiny\color{codegray},
  breaklines=true, captionpos=b, numbers=left, numbersep=5pt,
  keepspaces=true, showspaces=false, showstringspaces=false, showtabs=false,
  tabsize=2, frame=single, rulecolor=\color{gray!30},
  xleftmargin=10pt, xrightmargin=10pt,
}

\usepackage{enumitem}
\setlist{nosep}
\setlist[itemize,1]{label=--}
\usepackage{chngcntr}
\counterwithout{table}{chapter}
\counterwithout{figure}{chapter}
\renewcommand{\thetable}{\arabic{table}}
\renewcommand{\thefigure}{\arabic{figure}}
\usepackage{fancyhdr}
\pagestyle{plain}

\usepackage{titlesec}
\usepackage{tikz}
\usetikzlibrary{shapes.geometric, arrows, positioning, calc, er, fit}

\emergencystretch=5em
\tolerance=2000
\hyphenpenalty=300
\exhyphenpenalty=300
\titleformat{\chapter}{\bfseries\centering}{\thechapter}{1em}{}
\titleformat{\section}{\bfseries}{\thesection}{1em}{}
\titlespacing*{\section}{1.25cm}{*2}{6pt}
\titleformat{\subsection}{\bfseries}{\thesubsection}{1em}{}
\titlespacing*{\subsection}{1.25cm}{*2}{6pt}
\titleformat{\subsubsection}{\bfseries}{\thesubsubsection}{1em}{}


% babel переустанавливает \contentsname в \begin{document} — задаём через captions.
((* if project.lang == 'en' *))\addto\captionsenglish{\renewcommand{\contentsname}{Contents}}((* else *))\addto\captionsrussian{\renewcommand{\contentsname}{СОДЕРЖАНИЕ}}((* endif *))

\setcounter{tocdepth}{3}
\setcounter{secnumdepth}{3}

\hypersetup{unicode=true,breaklinks=true,bookmarksopen=true,bookmarksnumbered=true,colorlinks=true,linkcolor=black,urlcolor=blue,citecolor=blue}

% Значки для матриц сравнения (обзор аналогов/подходов).
\newcommand{\cmpplus}{\textbf{+}}
\newcommand{\cmpminus}{\textbf{--}}
\newcommand{\cmpplanned}{\textcolor{blue}{$\pm$}}

% Формат номеров в списке источников с точкой.
\makeatletter
\renewcommand\@biblabel[1]{#1.}
\makeatother

% ── Данные проекта: \hseFill, \signdateblank и все \hse-макросы ──
\input{../common/meta}

\begin{document}
\input{title}

\chapter*{Abstract}
\addcontentsline{toc}{chapter}{Abstract}
% GOST 7.32 §5.3: volume, keywords, then the abstract text — content and
% results, not just the topic.
The report has \pageref{LastPage}~pages, \ldots~figures, \ldots~tables, \ldots~references, \ldots~appendices.

\medskip
\hseFill{KEYWORD, KEYWORD, KEYWORD, KEYWORD, KEYWORD}

\medskip
The object of study is \hseFill{object}. The aim of the work is \hseFill{aim}.

\hseFill{what was done: methods, stages, experiments}.

\hseFill{results with numbers: metrics, characteristics of the solution}.

\hseFill{field of application, adoption, further development}.

\clearpage
\tableofcontents
\clearpage

\chapter*{Definitions and Abbreviations}
\addcontentsline{toc}{chapter}{Definitions and Abbreviations}
\begin{description}[leftmargin=0pt, labelindent=0pt, font=\bfseries, style=nextline]
  \item[\hseFill{Term}] \hseFill{definition}.
  \item[\hseFill{ABBR}] --- \hseFill{expansion}.
\end{description}

\clearpage
\chapter*{Introduction}
\addcontentsline{toc}{chapter}{Introduction}
\hseFill{relevance: the problem, who it affects and why existing solutions fall short}~\cite{gost732}.

\textbf{Object of study:} \hseFill{object}.
\textbf{Subject of study:} \hseFill{subject}.

\textbf{Aim:} \hseFill{the aim from the project proposal}.

The tasks are:
\begin{enumerate}[label=\arabic*), leftmargin=1.25cm]
  \item \hseFill{task 1};
  \item \hseFill{task 2};
  \item \hseFill{task 3}.
\end{enumerate}

((* if project.kind == 'research' *))
\textbf{Novelty:} \hseFill{what is new compared to known work}.
((* else *))
\textbf{Novelty and practical significance:} \hseFill{how the solution differs from analogues and where it is used}.
((* endif *))

\hseFill{report structure: what each chapter covers}.

((* if project.kind == 'research' *))
\clearpage
\chapter{Literature Review and Problem Statement}
\section{State of the art}
Review the work on the topic with citations.

\section{Comparison of approaches}
Compare the known approaches by relevant criteria (Table~\ref{tab:approaches}).

\begin{table}[H]
\caption{Comparison of existing approaches}
\label{tab:approaches}
\centering
\begin{tabular}{|>{\raggedright\arraybackslash}p{5cm}|c|c|c|}
\hline
\textbf{Criterion} & \textbf{Approach~A} & \textbf{Approach~B} & \textbf{Proposed} \\
\hline
\hseFill{criterion 1} & \cmpplus & \cmpminus & \cmpplus \\
\hline
\hseFill{criterion 2} & \cmpminus & \cmpplanned & \cmpplus \\
\hline
\end{tabular}
\end{table}

\section{Problem statement}
State the research question, the hypothesis and the formal problem.

\section{Chapter summary}
\hseFill{the open problem and the research task}.

\clearpage
\chapter{Theoretical Results}
\section{Model (method, algorithm)}
Describe the proposed model or method.

\section{Justification}
Give proofs of properties or complexity estimates.

\section{Chapter summary}
\hseFill{key results of the chapter}.

\clearpage
\chapter{Implementation and Experiment}
\section{Implementation}
Describe the prototype: architecture, technologies, limitations.

\section{Experimental setup}
Experiment plan, data and metrics.

\section{Chapter summary}
\hseFill{implementation and experiment outcomes}.

\clearpage
\chapter{Analysis of Results}
\section{Results}
Present the results (Table~\ref{tab:results}) against baselines.

\begin{table}[H]
\caption{Experimental results}
\label{tab:results}
\centering
\begin{tabular}{|>{\raggedright\arraybackslash}p{5cm}|c|c|}
\hline
\textbf{Method} & \textbf{Metric~1} & \textbf{Metric~2} \\
\hline
Baseline & \ldots & \ldots \\
\hline
Proposed & \ldots & \ldots \\
\hline
\end{tabular}
\end{table}

\section{Discussion and limitations}
Assess significance, threats to validity and applicability.

\section{Chapter summary}
\hseFill{the answer to the research question}.
((* else *))
\clearpage
\chapter{Domain Analysis and Problem Statement}
\section{Domain}
Describe the domain, the users and their tasks.

\section{Review and comparison of existing solutions}
Compare the analogues by functionality, efficiency and usability (Table~\ref{tab:analogues}).

\begin{table}[H]
\caption{Comparison of existing solutions}
\label{tab:analogues}
\centering
\begin{tabular}{|>{\raggedright\arraybackslash}p{5cm}|c|c|c|}
\hline
\textbf{Criterion} & \textbf{Analogue~A} & \textbf{Analogue~B} & \textbf{This work} \\
\hline
\hseFill{criterion 1} & \cmpplus & \cmpminus & \cmpplus \\
\hline
\hseFill{criterion 2} & \cmpminus & \cmpplanned & \cmpplus \\
\hline
\end{tabular}
\end{table}

\section{Problem statement and requirements}
State the problem, functional and non-functional requirements (details in
the requirements specification, Appendix~A).

\section{Chapter summary}
\hseFill{what is to be built and why}.

\clearpage
\chapter{Design}
\section{System architecture}
Describe the architecture, components and their interaction; justify the choices.

\section{Data model and interfaces}
Describe the data model, APIs and the user interface.

\section{Chapter summary}
\hseFill{key design decisions}.

\clearpage
\chapter{Implementation and Testing}
\section{Technologies and tools}
Languages, frameworks, infrastructure, development methodology and project
management tools.

\section{Implementation details}
Key decisions and size characteristics (components, lines of code, tests).

\section{Testing and evaluation}
Test method, results and comparison with analogues by the chosen metrics.

\section{Chapter summary}
\hseFill{implementation and testing outcomes}.
((* endif *))

\clearpage
\chapter*{Conclusion}
\addcontentsline{toc}{chapter}{Conclusion}
In this work:
\begin{enumerate}[label=\arabic*), leftmargin=1.25cm]
  \item \hseFill{outcome of task 1};
  \item \hseFill{outcome of task 2}.
\end{enumerate}
The aim has been achieved: \hseFill{how the expected results were confirmed}.
Future work: \hseFill{development in the master's thesis}.

\clearpage
\phantomsection
\addcontentsline{toc}{chapter}{References}
\begin{thebibliography}{99}
  \bibitem{gost732}
  GOST 7.32--2017. Report on Research Work. Structure and Rules of Presentation. Moscow: Standartinform, 2017.
\end{thebibliography}

\clearpage
\appendix
\titleformat{\chapter}[hang]
  {\normalfont\large\bfseries\centering}{Appendix~\thechapter}{1ex}{}
((* if project.kind == 'research' *))
\chapter{Source Code and Data}
Repository with the program and the experiment data: \hseFill{repository URL in the FCS GitLab}.
((* else *))
\chapter{Technical Documentation}
The technical documentation and the work plan are kept in the project
repository: \hseFill{repository URL in the FCS GitLab}. Access: \hseFill{how to get access}.
((* endif *))

\end{document}
